% Public copy: private contact details omitted at the owner’s request.
\documentclass[a4paper,10pt]{article}
\usepackage[margin=0.5in]{geometry}
\usepackage{enumitem}
\usepackage[hidelinks]{hyperref}
\usepackage{titlesec}
\usepackage{tabularx}

\titleformat{\section}{\large\bfseries}{}{0em}{}[\titlerule]
\setlist[itemize]{leftmargin=1.5em, topsep=2pt, itemsep=2pt}

\begin{document}

\begin{center}
{\LARGE \textbf{Michael Alexander Saxer}}\\[4pt]
B.Sc. Data Science \quad | \quad Electronics Engineer (VET)\\[4pt]
GitHub: \href{https://github.com/saxermi}{github.com/saxermi} \quad | \quad
LinkedIn: \href{https://linkedin.com/in/michaelsaxer}{linkedin.com/in/michaelsaxer} \quad | \quad
ORCID: \href{https://orcid.org/0009-0005-7907-8170}{0009-0005-7907-8170}
\end{center}


\vspace{0.5em}

\section*{Education}
\textbf{M.Sc. in \href{https://ethz.ch/en/studies/master/degree-programmes/engineering-sciences/computational-science-and-engineering.html}{Computational Science and Engineering}}, ETH Zürich \hfill \textit{09/2026 – 06/2029}\\
Physics specialization; individual 164 ECTS study plan.\
Focus on computational modeling, numerical methods, high-performance computing, and physics applications.\\
\textbf{B.Sc. in Data Science}, Zurich University of Applied Sciences (ZHAW) \hfill \textit{09/2022 – 06/2025}\\
\textbf{Thesis:} \textit{Optimization of Track Clustering for Primary Vertex Reconstruction in Proton-Proton Collisions at the CMS Experiment at CERN. DOI: \href{https://doi.org/10.21256/zhaw-34860}{10.21256/zhaw-34860}}\\
Specialization in machine learning and computational modeling; additional coursework in PDE numerics and engineering physics.\\
\textbf{Federal Vocational Baccalaureate, Engineering/Architecture/Life Sciences} \hfill \textit{08/2021 – 06/2022}\\
\textbf{Final Thesis}: \textit{Therapy of Prostate Cancer with Lutetium-177-PSMA-617.}\\
Bilingual program (English/German).


\vspace{0.5em}

\section*{Research Experience}

\textbf{Undergraduate Researcher}, Paul Scherrer Institute (PSI), High Energy Physics Group \hfill \textit{08/2024 – 06/2025}\\
Semester Project and Bachelor Thesis within the CMS Experiment at CERN, supervised by Dr. W. Erdmann (PSI / ETH Zurich).
Focused on primary-vertex reconstruction algorithms for the High-Luminosity LHC. As part of this work, I attended the Inverted CERN School of Computing (iCSC 2025) in Meyrin, with hands-on modules in HEP computing, machine learning, and software best practices. Both the semester project (6 ECTS) and bachelor thesis (12 ECTS) were graded 6/6.
\begin{itemize}[leftmargin=1.2em, topsep=2pt, itemsep=2pt]
    \item Evaluated the Deterministic Annealing in Blocks algorithm on HL-LHC simulation data and identified block-border effects as the main limitation for reconstruction quality under high-pileup conditions.
    \item Developed \textbf{CascadeDA}, a hybrid vertexing algorithm that combines block-wise pre-clustering with global deterministic annealing.
    \item Achieved 60\% faster clustering than the previous implementation on large-scale CMS tracking datasets while maintaining sub-10 µm z-resolution and a 99\% track-to-vertex assignment efficiency.
\end{itemize}



\vspace{0.5em}

\section*{Work Experience}

\textbf{Teaching Assistant}, Zurich University of Applied Sciences - School of Engineering \hfill \textit{08/2023 – 06/2024}\\
In the summer break, I taught mathematics tutorials in the ZHAW Pre-College program (technical vocational baccalaureate). During the semester, I supervised weekly tutorials for first-year engineering students in Calculus I–II and Linear Algebra I–II, and provided additional support in Python programming and study methods.
\vspace{0.5em}\\

\textbf{Apprenticeship in Electronics Engineering (VET Diploma)}, Siemens Switzerland AG \hfill \textit{08/2017 – 08/2021}\\
Training in electronic circuits, embedded systems, and software engineering. Developed projects in C, C++, and C\# and implemented building automation solutions.

\vspace{0.5em}
\section*{Publications and Presentations}

\textbf{L. Stuhlmann, M. A. Saxer, J. Fürst}, “\textit{Efficient and Reproducible Biomedical Question Answering Using Retrieval Augmented Generation},” \textit{IEEE Swiss Conference on Data Science (SDS)}, Zurich, 2025. DOI: \href{https://doi.org/10.1109/SDS66131.2025.00029}{10.1109/SDS66131.2025.00029}\\[4pt]
\textbf{Presentation:} \textit{Quality Assessment of Parallelized Primary Vertex Reconstruction at the CMS Experiment}, Tracker POG Meeting, CERN/Online, Feb 2025.(\href{https://indico.cern.ch/event/1498184/}{Link to Indico} (Access requires CERN Account))\\[4pt]
\textbf{Presentation:} \textit{Primary Vertex Reconstruction on GPU}, Tracker POG Meeting, CERN/Online, Dec 2025.(\href{https://indico.cern.ch/event/1616105/}{Link to Indico} (Access requires CERN Account))
\vspace{0.5em}
\newpage
\section*{Awards and Affiliations}

\textbf{Swiss Study Foundation - Mentoring Program} \hfill \textit{Since 2024}\\
Selected for academic excellence and intellectual potential among Swiss university students.\\[4pt]
\textbf{YMCA Scout Leader}, CEVI Altstetten-Albisrieden\hfill \textit{Since 2016}\\
Led activities and security operations for 1,500+ (2022 regional camp) and 500 participants (2024 national leaders assembly).\\[4pt]
\textbf{President}, FDP.The Liberals Zurich District 9 \hfill \textit{Since 2026}\\
Representing the District 9 branch, VP since 2022; ran for communal (2022/2026) and cantonal (2023) elections. First in line to fill a vacancy in Zurich’s City Parliament.

\vspace{0.5em}

\section*{Projects}
\href{https://github.com/Saxermi/hana_pde_modeling_fem/blob/main/Hana_projekt.pdf}{\textbf{Modeling Charged Particle Motion in Magnetic Fields with FEM and PINNs}} \hfill Grade: \textbf{6/6}\\
Modeled magnetic fields with the finite element method in NGSolve and simulated charged particle trajectories in Python using a fourth-order Runge-Kutta scheme, then implemented and evaluated a physics-informed neural network (PINN) with PyTorch for the Lorentz force equation.
\vspace{3mm}\\
\href{https://github.com/slinusc/speaker_identification_evaluation/blob/main/Evaluating_the_Effectiveness_of_Transformer_Layers.pdf}{\textbf{Layer-wise Analysis of Transformer Models for Speaker Identification}} \hfill Grade: \textbf{6/6}\\
Evaluated Wav2Vec 2.0, XLS-R, and Whisper for speaker identification by fine-tuning and analyzing layer-wise representations (SVCCA, k-means, t-SNE), showing different depth-dependent behavior of speaker features and identifying the optimal number of transformer layers for each model.


\vspace{0.5em}




\section*{Skills}

\begin{tabularx}{\linewidth}{@{}lX@{}}
\textbf{Programming:} & Python, C, C++, C\#, R \\[2pt]
\textbf{ML \& Data:} & PyTorch, scikit-learn, RAG, NLP, RL, Computer Vision \\[2pt]
\textbf{Scientific Computing:} & FEM (NGSolve), deterministic annealing, SLURM, ROOT, CMSSW \\[2pt]
\textbf{Tools:} & Git, Linux, Docker, \LaTeX, Flask \\[2pt]
\textbf{Languages:} & German (native), English (fluent), French (basic)
\end{tabularx}

\end{document}
